\documentclass[10pt]{article}
\usepackage[margin=20mm]{geometry}
\usepackage{amsmath,amssymb,fontspec,unicode-math,xurl,hyperref}
\IfFontExistsTF{DejaVu Serif}{\setmainfont{DejaVu Serif}}{\setmainfont{Latin Modern Roman}}
\setmathfont{latinmodern-math.otf}
\setlength{\parindent}{0pt}\setlength{\parskip}{6pt}\sloppy
\begin{document}
ELVIA Scientific Verification Report 0071
Position: Current Model Licensing Practices are Dragging Us into a Quagmire of Legal Noncompliance
Paper I25-40180 | Supported outcome: Minor
Review date: 24 October 2026 | Human scientific decision pending\par
\# 1. Identity, version and review result\par
Reviewed source: Moming Duan, Mingzhe Du, Rui Zhao, Mengying Wang, Yinghui Wu, Nigel Shadbolt and Bingsheng He. ICML 2025, PMLR 267:81276-81291. The exact 16-page publisher-linked PDF was supplied by the controller and its SHA-256 was independently verified against controller\_acquisition\_receipt.json. Page references below are 1-based PDF pages, equal to the printed internal page numbers, not the PMLR volume pagination.\par
Supported outcome: three bounded substantive issues (Minor): a mistaken Creative Commons definition, an unjustified conversion of method composition into method-specific risk, and an inconsistent Table 3 percentage. Compatibility exceptions and the La2E/La3E license-tag rule remain scoped Needs Evidence items. No executed classifier or corrected corpus prevalence is claimed.\par
The review read all main sections and Appendices A-E, checked the central quantitative definitions against Table 3, inspected primary license texts and current author project information, and independently recalculated the relevant arithmetic. PDF pages 1 and 13 were visually inspected. No author software or models were executed. Scientific acceptance and author adoption remain pending.\par
\# 2. Definitions, dependencies and coverage\par
The paper crawls a popularity-selected Hugging Face graph in December 2024: 151,302 model nodes and 146,818 relationships from the top 8,000 downloaded models and their trees (Section 2, p.3). Its NCM is the number of distinct flagged models; CTM is the number deemed eligible for a given noncompliance case (Appendix C, p.13). Method columns summarize adjacent dependency types among flagged models. Multiple paths are allowed, so those columns can sum to more than NCM.\par
The position concerns model licensing standards. Its actual technical support includes a license-clarity rubric, metadata-based conflict rules, descriptive counts, and a proposed license family. No experimental efficacy requirement is imposed on its normative recommendation. Copyrightability, license scope, source provenance, and enforceability are not treated as interchangeable facts.\par
Coverage: Sections 1-2 and Figures 1-3 were checked for identity, definitions and sampling scope; Section 3 and Tables 1-3 for stated rules and numerical interpretation; Sections 4-6 for alternative positions and dependence; Appendices B-C for classification and counts; Appendix D for metadata limitations; Appendix E for the stated ModelGo sheet. Table 1’s entire legal scoring matrix and individual repository classifications were not independently recoded. The raw December 2024 graph and the exact generating classifier were not obtained from the inspected publisher/project links. No empirical prevalence was reproduced.\par
\# 3. Problem 1 - A conditional CC clause is turned into a definition\par
Source claim and suspect step: p.1, Introduction, License Mismatch paragraph, represents Licensed Material as limited to music, performance or sound recordings and uses that premise to motivate ambiguity when CC is applied to model weights.\par
Complete source check: CC BY 4.0 Section 1(a), and BY-SA 4.0 Section 1(a), contain a conditional rule about synchronizing musical material with a moving image. The actual general definition of Licensed Material is in BY Section 1(f) and BY-SA Section 1(h), and encompasses artistic or literary works, databases and other licensed material. Removing the condition changes the logical meaning. Thus this particular scope-based justification is false. [R2-R3]\par
Alternative reading and impact: ambiguity about model copyright, model-specific governance, software patents or source-code obligations can still be discussed. Those questions do not make the quoted definition accurate. The paper’s broader rationale has other premises and its Table 1 rubric does not explicitly compute scores from this mistaken restriction. The demonstrated damage is to this concrete supporting argument, so Minor is supported; a change to the entire mismatch count is not established.\par
Scoped correction: replace the claimed definition with the actual general definition and discuss model-specific gaps directly. This restores the local source attribution. It does not validate every Table 1 entry or establish the legal status of model weights.\par
\# 4. Problem 2 - Within-case composition does not determine conditional risk\par
Source claim: Appendix C, pp.13-14, introduces the method columns as an analysis of which reuse methods are more prone to noncompliance, then concludes that certain methods are more likely to lead to it. In the copyleft-violation row, 296 finetune, 32 merge and 37 quantization flags sum to NCM=365; CTM=949. The displayed 81.1\%, 8.8\% and 10.1\% are fractions within the flagged set.\par
The suspect inference reverses the conditioning. Counts of flagged methods estimate composition among flagged models, not the probability of a flag among all eligible models using each method. The latter requires method-specific eligible denominators. Neither the row nor the global graph’s method proportions supplies those denominators for this particular eligible subset.\par
Complete logical counterexample: assign the 949 eligible models to disjoint method groups of 850 finetune, 32 merge and 67 quantization models, with 296, 32 and 37 flags respectively. Every flag count is nonnegative and at most its denominator, the flags total 365, the eligible models total 949, and the row’s three reported composition percentages are preserved. Disjoint labels are an allowed special case: the paper permits multiple paths but does not require them. Yet conditional flag rates are 34.8235\%, 100\% and 55.2239\%, respectively. Finetune has the largest number of flags and the lowest rate. These are reviewer-constructed values demonstrating non-identifiability, not recovered corpus data.\par
Counterevidence and impact: global finetune prevalence is reported in Section 2, but it concerns the full graph, not the copyleft-eligible subset. Multi-path counting creates an additional denominator choice rather than resolving this one. This invalidates the method-specific comparative-risk inference from the supplied table. It does not change the reported marginal counts or the rule used to detect a case. Minor is assigned because this is a secondary stratified interpretation with no demonstrated dependence of the principal conflict detector or overall prevalence on it; its appendix location is not the reason.\par
Scoped correction: call the columns the composition of flagged cases, or provide and audit eligible denominators for every method before making a comparative-risk statement. A corrected empirical ranking requires the raw graph; none is invented here. The finite calculation is replayable in independent\_checks.py.\par
\# 5. Problem 3 - Table 3’s copyleft-style percentage is internally inconsistent\par
Source: p.13, Table 3, Copyleft-style Terms row: NCM=26,679, CTM=41,019 and 63.5\%; p.7 repeats 63.5\%. By the paper’s declared NCM/CTM definition, 100 x 26,679 / 41,019 = 65.04059094565933\%, which rounds to 65.0\%, a difference of 1.54059 percentage points. This was checked from the rendered original table and with independent scalar arithmetic.\par
The three entries cannot all be correct simultaneously. The evidence establishes a local inconsistency, not which original integer is wrong. If both printed counts are intended, 65.0\% is the precise one-decimal repair. If another denominator generated 63.5\%, that denominator and its eligibility definition must be supplied. Minor: this changes one reported rate and its repetition, without demonstrating reversal of the broader descriptive observation that the flagged fraction is large. Raw-data repair remains unverified.\par
\# 6. Unresolved checks, counterevidence and known limitations\par
Compatibility (Needs Evidence for empirical impact). Section 3.2, p.6, and Appendix B use broad same-license wording; Appendix C, p.13, similarly generalizes OSS copyleft restrictions as CC conflicts. However, p.7 expressly recognizes BY-SA 4.0 to GPLv3 compatibility, and p.12 explicitly permits a compatible license. The official CC compatibility list confirms the exception. A literal rule that rejects every different identifier would be too broad, but the paper contains a defensible exception-aware reading. Without the exact classifier, no confirmed miscount among the 365 copyleft or 76 CCf cases follows. [R3-R4]\par
La2E/La3E license-tag inference (Needs Evidence). Section 3.3, p.7, flags a Llama2/3 license label when the corresponding base is absent. Meta’s relevant Section 1.b.v concerns using the corresponding Meta materials or outputs to improve another LLM. A tag alone does not establish this antecedent. Conversely, an absent edge does not establish absence of actual use. A constructed independently trained model with an inappropriate Llama tag would be a metadata/license-scope problem; it would not by itself prove a violation of the cited use clause. Because that construction may exploit the very input errors acknowledged by the authors, it is not promoted here to a demonstrated Major defect under verified operative assumptions. The direct dependence on the 1,000 La2E-license and 966 La3E-license flags warrants targeted source-level adjudication. These counts are not added as if disjoint. [R5-R6]\par
Known-issue reproduction: Appendix D, pp.15-16, already discloses missing/inaccurate dependencies and license metadata, including 276 La3E and 18 La2E candidate identifiers suggestive of false positives. Their disclosure is preserved and attributed to the authors. This review does not identify those models, adjudicate them, subtract them from totals, or relabel that disclosure as a newly discovered error.\par
Later-version check: the current ModelGo website reports changes in 2026 and a withdrawal of its share-alike variant from OSI review. This is not treated as a retroactive defect or a verified repair of the paper’s November 2024 ModelGo 2.0 comparison. A version-matched historical license text and complete scoring audit remain outstanding. The publisher page was inspected; OpenReview’s forum presented a browser challenge and its public note API returned HTTP 403. No accessible paper erratum was identified in the inspected sources; this is not an exhaustive guarantee of absence. [R1,R7]\par
\# 7. Dependence, repair status and conclusion\par
The supported errors affect one foundational example, one secondary comparative-risk inference and one table rate. They do not establish that model-license conflicts never occur, that all reported counts are wrong, or that a standardized license proposal cannot be useful. Neither ease of repair nor survival of the position determines severity: the bounded dependencies above do.\par
The unresolved classifier questions could affect central quantitative evidence, but the missing implementation and alternative readings prevent a demonstrated severity upgrade. Appropriate follow-up is to inspect the frozen graph, versioned classifier, compatibility exceptions and per-case adjudications; this requires source evidence, not model training.\par
Item | Judgment | Repair verification | Remaining work\par
P1 CC definition | Minor | Local wording correction supported | Audit any dependent scoring\par
P2 method-risk inference | Minor | Logical reversal demonstrated | Obtain eligible denominators\par
P3 printed percentage | Minor | 65.04059\% if counts fixed | Resolve original count source\par
CC compatibility / Llama tags | Needs Evidence | No empirical repair claimed | Inspect exact classifier and cases\par
Human decision / author adoption | Pending | Not requested or supplied | Controller acceptance and human review\par
\# 8. Consulted sources\par
[R1] Duan et al., ICML 2025, exact publisher PDF and metadata: https://proceedings.mlr.press/v267/duan25d.html ; source PDF: https://raw.githubusercontent.com/mlresearch/v267/main/assets/duan25d/duan25d.pdf .\par
[R2] Creative Commons, Attribution 4.0 International legal code, Sections 1(a), 1(f), 2 and 3: https://creativecommons.org/licenses/by/4.0/legalcode.en .\par
[R3] Creative Commons, Attribution-ShareAlike 4.0 International legal code, Sections 1(a), 1(h), 3(b): https://creativecommons.org/licenses/by-sa/4.0/legalcode.en .\par
[R4] Creative Commons, Compatible Licenses, BY-SA versions 3.0 and 4.0, retrieved 24 October 2026: https://creativecommons.org/compatible-licenses/ .\par
[R5] Meta, Llama 2 Community License Agreement, release date 18 July 2023, definitions and Section 1.b.v: https://raw.githubusercontent.com/meta-llama/llama/main/LICENSE .\par
[R6] Meta, Meta Llama 3 Community License Agreement, release date 18 April 2024, definitions and Section 1.b.v: https://raw.githubusercontent.com/meta-llama/llama3/main/LICENSE .\par
[R7] ModelGo project website, current update history retrieved 24 October 2026: https://www.modelgo.li/ . This current website is a later-version source, not the frozen 2024 license text.\par
\# Appendix A. Evidence, replay, failures and roles\par
Primary PDF SHA-256: `580676cd424e45a21a6c468d11dedc9776f0fbfed480bdf88d9e098e3b58cdf8`. Source preservation and extraction are recorded in `verification/0071/source\_verification.json`. Exact supplementary HTTP bytes, hashes, timestamps and failed requests are in support\_acquisition.json; all downloaded support is under `sources/0071/`.\par
Reviewer-owned preparation: python `verification/0071/prepare\_source.py`. It verifies the hash, uses the existing Poppler pdftotext -layout -enc UTF-8 command, and saves pdfinfo. Finite replay: python `verification/0071/independent\_checks.py`. Output: `verification/0071/independent\_checks.json`. The scripts are separate support files and were actually executed. They do not import author code or call models.\par
Tool failures: Lumen’s PDF read\_file operation failed with Cannot read properties of undefined (reading ’toLowerCase’); existing Poppler supplied extraction and rendering. Default and bigbro Python environments did not expose fitz, pypdf or reportlab. No software was installed. The predecessor PDF was supplied through the document renderer. This replacement PDF was compiled offline from the accompanying TeX during acceptance; controller promotion remains pending.\par
Roles: ChatGPT performed source inspection, independent calculations, report preparation and visual checks. A second AI reviewer independently checked license interpretation and counterevidence within this same paper. Paper 01/02 evidence workflows were used; unavailable Nature-reviewer and paper-lookup skills were not claimed. No human verification or author endorsement is asserted.\par
Task: B001 / [internal task] / attempt-001. Batch token: [internal identifier omitted]. The exact task packet was read through Lumen. No queue, controller or final archive file was changed. Findings, coverage and provenance are supplied separately.\par
\end{document}
